\ssscp{Distribution of words for ‘parent of child’s spouse’}{15-04-03-02}
Looking at social structure,
\citet[204--210]{Fox2015} has shown that the line of demarcation for
where various terms meaning ‘child’s spouse’s parent’ are found
and not found is remarkably close to the Blust line in some places
as shown in \frf{15-05}.
But note that the feature is present in some SHWNG languages
(SHWNG has been classified within CEMP),
so the line delineating terms for ‘child’s spouse’s parent’ falls
to the east of Halmahera,
rather than the west.
Terms for such a relationship are pervasive west of the line in \frf{15-05}.
They are not found east of the line.
As \citet{Fox2015} points out,
the social structures of ethnic groups speaking Austronesian
languages within our target region are significantly different
in many ways from the social structures west and north of our target region.

\begin{figure}[h]\centering
	\includegraphics[width=\columnwidth]{besan.pdf}
	\fcp{Presence and absence of ‘parent of child’s spouse’}{15-05}
\end{figure}